\section{Future Work}

Several directions emerge naturally from the present work. A first is to
further validate and refine the assumed distribution for the error between the
saturated catalyst response and the actual system response. In particular,
fitting a more general Weibull distribution to the error data, using an
expectation maximization (EM) algorithm for parameter estimation, could provide
a more accurate statistical description and tighter confidence bounds for
saturation detection and saturated model parameters based aging detector. The
asymptotic approximations underlying the Wald test could be supplemented with
finite-sample corrections. A broader validation campaign using a larger fleet
of trucks with controlled and documented aging levels (e.g., accelerated aging
on a dynamometer followed by road testing) would enable rigorous calibration of
the detector threshold via the Neyman-Pearson criterion. Such a study would
also quantify how the non-centrality parameter $\lambda$ scales with the true
degree of catalyst degradation, establishing the minimum detectable aging level
for a given confidence requirement.

In parallel, the impact of the proposed nonlinear dynamic model on urea-dosing
control should be investigated. By explicitly exploiting the model structure,
future studies can assess whether more informed control policies improve
de-$NO_x$ efficiency, for example by deliberately maximizing the duration of
safely maintained catalyst saturation while respecting ammonia slip and
hardware constraints.

From a modelling standpoint, the Self-Excited Threshold Autoregressive (SETAR)
framework offers a promising avenue for capturing regime-switching behavior in
SCR-ASC dynamics. Alternative formulations for the temperature dependence of
reaction and transport parameters should be explored, with the goal of
achieving a wider temperature validity range and uniquely
identifiable parameter estimates across operating conditions. A natural model
extension is to incorporate additional reaction pathways, in particular, the
fast and slow SCR reactions and explicit ASC oxidation dynamics, into the
switched nonlinear framework. Where inter-brick measurements (between SCR and
ASC sections) or tailpipe ammonia sensors become available, these additional
signals could decouple the lumped rate constants and improve model fidelity,
especially at lower temperatures where the $100\%$ nitrogen selectivity
assumption breaks down. On the sensing side, explicit estimation and correction
of the $NO_x$ sensor cross-sensitivity to ammonia would sharpen the separation
between aged and degreened parameter estimates. 

Finally, integration of the proposed detector into an engine control unit (ECU)
prototype, with real-time parameter estimation running alongside the existing
urea-dosing controller, would provide a practical assessment of computational
feasibility, memory requirements, and latency. 